% =====================================================================
%  tr01b —— §3 讲义（原文 pp.10--11）
%  归属：tr-01。WP 辛几何、定理 2（p.10）、增广 Teichmüller 空间、映射类群、
%  Bers 纤维空间与万有曲线；图 4（p.11）。本块无编号公式。
% =====================================================================

% [原文 p.10]
Weil--Petersson（WP）度量是 K\"ahler 的。辛几何始于辛形式
$\omega=2\omega_{\WP}$ \emph{K\"ahler}，以及基本的扭转--长度对偶
\[
\omega(\ ,t_\alpha)=\dd\ell_\alpha.
\]
由 Lie 导数方程 $L_X\omega(\ ,\ )=\dd\omega(X,\ ,\ )+\dd(\omega(X,\ ))$
可得，形式 $\omega$ 在所有扭转流（twist flows）下不变。由此，测地长度函数是
FN 无穷小扭转的 Hamilton 势。对称性的论证表明，$\ell$ 与 $\tau$ 为 $\omega$
提供作用--角坐标。

% [Tutor Note: p.10 第一句"辛形式 $\omega=2\omega_{\mathrm{WP}}$ Kähler"里，紧跟在
%  公式后面的斜体词 Kähler 在原文中就是一个独立的斜体词（已裁 600 dpi 复核），
%  不是下标的一部分，应读作"即 Kähler 形式"；本译本保留原排版。真正要紧的是那个
%  因子 $2$：本讲义的 $\omega$ 取的是 $\mathrm{WP}$ 度量的 Kähler 形式的两倍。
%  核对 p.12 定理 3 的 $2\pi^2\kappa_1=\omega$ 与 p.13 定理 4 的 $C_\alpha$ 时，
%  必须先把 $\omega$、$\psi$ 的归一化与本约定对齐，否则与别的教材逐项比对会整体
%  差一个常数因子（$\omega$ 换一代数，$\omega^{3g-3+n-|\alpha|}$ 就整体变一次）。]

\begin{thm}{2} (W)，[Wlp10]。$\dd\ell\wedge\dd\tau$ 公式。WP 辛形式为
\[
\omega=\sum_{j=1}^{3g-3+n}\dd\ell_j\wedge\dd\tau_j.
\]
\end{thm}

边界空间（frontier spaces）被并到 $\T$ 上，这对应于允许 $\ell_j=0$，此时 FN
角 $\theta_j=2\pi\tau_j/\ell_j$ 无定义（在极坐标中，原点处的角度无定义）。长度
消失所描述的是这样的裤子：其相应的边界由穿刺点表示——方程 $\ell_\alpha=0$
描述的是 $\alpha$ 由一对穿刺点表示的双曲结构。对指标子集
$J\subset\{1,\dots,3g-3+n\}$，$J$-零化层为
\[
\cS(J)=\{R\ \text{退化}\mid \ell_j(R)=0\ \text{当且仅当}\ j\in J\}.
\]
每个零化层都是低维 Teichm\"uller 空间之积。$\T\cup\cS(J)$ 中邻域的一组基，
由参数 $(\ell_k,\theta_k,\ell_j)_{k\notin J,\ j\in J}$ 给出。增广
Teichm\"uller 空间是分层空间
\[
\Tbar=\T\cup_{\text{裤子分解 }\cP}\ \cup_{J\subset\cP}\cS(J).
\]

% [Tutor Note: p.10 这里有一笔维数账：$J$-零化层 $\cS(J)$ 的复维数是
%  $3g-3+n-|J|$，因为每令一个 $\ell_j=0$ 就少掉一个长度参数；而 FN 角
%  $\theta_j=2\pi\tau_j/\ell_j$ 在 $\ell_j=0$ 处失效，所以邻域参数只能用
%  $(\ell_k,\theta_k,\ell_j)_{k\notin J,\ j\in J}$ 这套混合坐标来写。原文紧接的
%  断言"$\Tbar$ 从不局部紧"没有给论证，它依赖于各层在边界处彼此接触的方式，严格
%  讨论见 [Wlp10] 关于增广空间与 Chabauty 拓扑的章节，本译注不展开。另注意
%  "每个层都是辛的"可由定理 2 直接限制到 $\cS(J)$ 上得到：$\dd\ell_j$ 在
%  $\ell_j=0$ 的各层上为零，剩下的 $\dd\ell_k\wedge\dd\tau_k$ 之和仍非退化。]

空间 $\Tbar$ 也可以这样描述：$\pi_1(F)$ 到 $PSL(2;\R)$ 的离散、忠实、保型的
表示，模去 $PSL(2;\R)$ 共轭，在 Chabauty 拓扑下的闭包。增广构造对 $\Tg$、
$\Tg(L)$ 与 $\Tgn$ 都有效。$\Tbar$ 是 Baily--Borel 型部分紧化。$\Tbar$ 从不
局部紧。$\dd\ell\wedge\dd\tau$ 公式给出了辛结构到增广 Teichm\"uller 空间
$\Tbar$ 的延拓。每个层（strata）都是辛的。

映射类群（MCG）$\mathrm{Homeo}^+(F)/\mathrm{Homeo}_0(F)$ 通过前复合作用在
标记上，从而作用在 $\T$ 上。对 $\mathrm{Homeo}^+(F)$，我们考虑保型（边界点、
边界曲线）、保持边界标号、保向的同胚。$\mathrm{Homeo}_0(F)$ 是相对边界（把边界作为集合）
同伦于恒等的元素构成的正规子群。一个 Dehn 扭转是这样一个同胚：它在一条简单闭
曲线的管状邻域之外是恒等，在同伦意义下非平凡，并且把管状邻域的一个边界相对于
另一个边界旋转。Dehn 扭转类生成 MCG。MCG 在 $\T$ 上真不连续地作用，且对
$\Tg$ 与 $\Tgn$ 通过双全纯映射作用。除有限多个拓扑型之外，该作用是有效的。
有限的 MCG 子群作用带不动点。MCG 作用在分层空间 $\Tbar$ 上。Bers 观察到存在
只依赖拓扑型的常数 $b_{g,n}$，

% [原文 p.11]
使得亏格 $g$、$n$ 个穿刺点的双曲曲面有一条缝合长度（seam lengths）不超过
$b_{g,n}$ 的裤子分解。由此，FN 坐标中的区域
$\{\ell_j\le b_{g,n},\ 0<\tau_j\le\ell_j\}$ 是一个粗略基本集——每条 MCG
轨道与该区域相交有界正数次。$\Tbar/\MCG$ 是紧的实解析 orbifold；
$\Tbar_g/\MCG$ 与 $\Tbar_{g,n}/\MCG$ 在拓扑上是 $\Mg$ 与 $\Mgn$ 的
Deligne--Mumford 稳定曲线紧化。Bers 关于基本集的观察与 $\dd\ell\wedge\dd\tau$
公式相结合，给出 $\M$ 的 WP 体积有限。

Bers 纤维空间 $\cC$（具体地说是 $\cC_g$ 与 $\cC_{g,n}$）是 $\T$ 上的复圆盘
全纯丛，它在 $\{(f,R)\}$ 上的纤维是万有覆盖 $\wt{R}$。纤维上的一点可以看作
一个穿刺点，并确定基本群中从基点出发的一条曲线——因此
$\cC_g\approx\T_{g,1}$，$\cC_{g,n}\approx\T_{g,n+1}$。把 $\MCG(F)$ 用基本群
$\pi_1(F)$ 扩张所得的群 $\MCG_\cC$，真不连续且全纯地作用在 $\cC$ 上。群
$\MCG_\cC$ 同构于 $\MCG_{g,n+1}$。对从 $\MCG_\cC$ 到 $\MCG(F)$ 的满同态，
前一个群等变地作用在 $\cC$ 到 $\T$ 的纤维化上。由此得到的映射
$\pi:\cC/\MCG\to\T/\MCG$ 描述了一个 orbifold 丛，即万有曲线，其纤维是
orbifold 纤维——纤维是模去其全自同构群的 Riemann 曲面。流形的有限局部覆盖
与其商可以这样来描述：从具有局部极大对称性的曲面出发，用曲面的典范（极值、
调和）映射来引入丛的局部平凡化。

% [Tutor Note: p.11 的 $\cC_g\approx\T_{g,1}$、$\cC_{g,n}\approx\T_{g,n+1}$ 是
%  维数锚点：Bers 纤维空间 $\cC_g$ 是 $\T_g$ 上的复圆盘丛，复维数
%  $(3g-3)+1=3g-2$，正是 $\T_{g,1}$ 的复维数 $3g-3+1$；同理 $\cC_{g,n}$ 的复
%  维数 $(3g-3+n)+1=3g-3+(n+1)$ 对应 $\T_{g,n+1}$。原文把扩张群与商都记作
%  $\MCG_\cC$，读的时候要按上下文区分"作用在 $\cC$ 上的扩张群"与"商里的群"。]

增广构造应用于 Bers 纤维空间，得到 $\overline{\cC}$。

% [原文 p.11] 图 4
\begin{figure}[htbp]
\centering
\includegraphics[width=0.55\textwidth]{fig04.png}
\figtitle{4}{万有曲线 $\overline{\cC}/\MCG_{\overline{\cC}}$ 在 $\Tbar/\MCG(F)$ 上的一个穿刺截面。}
\end{figure}

$\MCG_\cC$ 作用在分层空间 $\overline{\cC}$ 上。增广商
$\overline{\cC}/\MCG$ 是一个 orbifold，并且在 $\Mbar$ 上几乎是一个 orbifold
丛——在 Riemann 曲面的一个节点（一对穿刺点）处，纤维变为竖直的——纤维化的
局部模型是投影 $\{(z,w)\}\to\{t=zw\}$ 在原点处的芽，即复双曲线族。该几乎
orbifold 丛 $\overline{\cC}/\MCG\to\Mbar$ 的纤维的转动（turning）由 $\Mbar$
上的线丛来度量。纤维的切 $\C$-线族 $(\mathrm{Ker}\,\dd\pi)$ 是沿光滑 Riemann
曲面的切丛，
